\section{From overlapping key distributions to four holes}\label{sec:collision}

The deterministic construction becomes useful when its keys collide.
The principal cost is equality of ambient vectors; the remaining
primal-channel scalar entries have a smaller cost.  We give a
criterion that retains this distinction, including for cells selected
using the parameters on the opposite side.

\subsection{Queries and their reference measure}

A numerical template fixes the shape of all lists in a scalar
recipe: positions, endpoints, tags, labels, flavors and desired
$q$ bits, with matching tags and $q$ bits across the two sides.
It also fixes each atom's \emph{numerical unary status}: its tester
summary, role $a(\mathrm{atom})$, basis values of
$T(\mathrm{atom},\cdot)|_{C_i}$, and $p_i$ bit.
It fixes the within-endpoint $L,R$ bit prescriptions of
\Cref{lem:small-tables}. Its labels are distinct at every endpoint.
For a fixed pair of leaves it may specify a small table and two
separate orientation filters, including its admissibility filters.
Only choices for which equal accepted keys satisfy the hypotheses
of \Cref{thm:gradient-realization} will be used.

The query at a unit draws fresh independent point parameters at
all its positions and evaluates their keys.  Parameters may be
uniform on the retained flavor coordinates, or normalized
conditional on the prescribed, possible $q$ bit at each position.
This latter conditioning depends only on point parameters and
has a fixed positive probability whenever it is possible.  In
particular the two units' parameter draws remain independent.
No normalization is made for the other status, binary, Gram or
orientation acceptance conditions.  Every resulting accepted
measure is a subprobability measure.

Let $k$ be the number of ambient vector slots on one template
side, counting components and signs.  There are at most four
lists of seven atoms, each with $2(g-1)$ slots, so
\begin{equation}\label{eq:key-rank-budget}
 k\le 56(g-1)=k_{\max}.
\end{equation}
Let $\cK_n$ be the set of global tuples of these $k$ vectors
satisfying the prescribed paired diagonal bit at every position
and component, and having every off-position plus/minus inner
product within a component equal to zero.  Write $\nu_n$ for
the uniform probability measure on $\cK_n$.

The set $\cK_n$ has density bounded below by a positive constant
in the space of all $k$ independent uniform vectors, and
\begin{equation}\label{eq:key-space-size}
 |\cK_n|\le 2^{kN}.
\end{equation}
Indeed the number of slots is fixed.  With probability tending
to one the plus vectors in each component are independent;
given them, each minus vector fulfills a fixed list of
independent linear equations with the prescribed probability
$2^{-r}$, where $r$ is the number of equations on it.  Their
total number is fixed.  This supplies a constant lower bound,
including all diagonal and off-position requirements.
Queries are accepted only when their keys belong to $\cK_n$.

For a pair of leaves, let $d_A,d_B$ be the densities relative
to $\nu_n$ of the separate accepted key measures, integrating
over their normalized leaf orientation laws and parameter draws.
The densities include their orientation filters and all query
conditions but are not renormalized for success.  Thus
$\int d_A\,d\nu_n,\int d_B\,d\nu_n\le1$.

\begin{proposition}[Collision criterion]\label{prop:collision}
Let $\sigma$ be a unit law in \Cref{thm:raw-supersaturation},
and let $\widetilde\sigma$ be a restriction of it of relative
mass at least $2^{-o(N)}$, decomposed into leaves satisfying
\eqref{eq:final-leaf-minentropy}.  Mix the leaves with their
actual masses in $\widetilde\sigma$.
Suppose accepted numerical queries as above, with their
leaf-pair choices of table and separate filters, satisfy
\begin{equation}\label{eq:overlap-criterion}
 \E_{\text{independent leaf pair}}
       \int d_A d_B\,d\nu_n\ge 2^{-o(N)}.
\end{equation}
Then two independent units from the original $\sigma$ have
all four cross holes with probability at least
\begin{equation}\label{eq:collision-conclusion}
 2^{-(k_{\max}+.03)N-o(N)}.
\end{equation}
It suffices to obtain \eqref{eq:overlap-criterion} after skipping
some leaf pairs, or after summing over a bounded number of
numerical options.  Positive subsequential constants and
inverse-polynomial overlap bounds both suffice.  In particular,
this conclusion contradicts a sequence violating
\Cref{thm:raw-supersaturation}.
\end{proposition}
\begin{proof}
We first work with one shape of vector slots and one option per
leaf pair.  A bounded number of shapes or options will be
handled at the end.

\paragraph{The numerical records.}
Fix two leaves, a key $q\in\cK_n$, and point parameter values
$a,b$ on the two sides.  Freeze the images of all pin and key
directions.  The pin images are known from the leaves and the
key images are the entries of $q$.  At each unit record the
pairings of all its nominal columns with these opposite frozen
images.  Record any numerical statuses in
\Cref{eq:recipe-scalars,eq:recipe-gradients} not already fixed
by the query conditions.  For fixed parameters these records,
together with acceptance, partition the unit's orientations
into unary cells.  A partition is allowed to depend on the
opposite parameters; it does not depend on the opposite
orientation once its frozen images have been specified.

There are at most
\begin{equation}\label{eq:record-count}
 L_n=2^{Cn}
\end{equation}
cells on either side, for a constant $C$ fixed before $M_0$.
To verify the count, there are $O(\dim\cB+h)=O(n)$ columns
and only boundedly many pin and key vectors to pair against.
Nominal projected-pin bases use $O(n)$ coefficients.  The
basis tensors of $C_i$ are encoded in their bounded projected
pin bases, rather than as full matrices on $\cB$; this uses
boundedly many further coordinates.  Parameter values require
$O(n)$ bits and the remaining table and status information has
bounded size.  The ambient pin and key images are already
fixed and are not charged again as new records.

The numerical records determine every frozen cross entry.
Together with the fixed algebraic data, leaves, parameters,
and table, they determine a choice of the abstract solution
in \Cref{thm:gradient-realization}.  One can make this choice
deterministic by taking the first solution in fixed coordinate
orders on the finite spaces.  Its unknown entries need not
be encoded as additional records.  Any further orientation
acceptance test is only a predicate; its description is not
an additional input to this solver.

\paragraph{Pruning at a fixed parameter pair.}
Put
\begin{equation}\label{eq:cell-threshold}
 \tau_n=2^{-(k+.01)N}.
\end{equation}
For the fixed parameter pair $(a,b)$ and key $q$, let
$\alpha_a^b(c;q)$ be the mass of a first-side accepted
orientation in record cell $c$, measured in its normalized
leaf law.  Define $\beta_b^a(d;q)$ similarly.  Denote their
sums over records by $\alpha_a(q),\beta_b(q)$; these sums
do not depend on the opposite parameter, since it only
refines the record partition.  Their sums over all keys
are at most one.

Discard for this fixed pair only cells of mass less than
$\tau_n$.  The product mass lost from small first-side cells
at key $q$ is at most
\[
 L_n\tau_n\,\beta_b(q).
\]
Sum over keys and average the independent parameter weights.
The loss in collision probability is at most $L_n\tau_n$,
since $\sum_q\beta_b(q)\le1$ for every $b$.
The reciprocal loss is at most the same amount.
Because $\nu_n$ is uniform, the overlap integral is
$|\cK_n|$ times the probability of an accepted key collision.
Thus the total overlap loss, for any fixed leaf pair and
also after averaging leaf pairs, is at most
\begin{equation}\label{eq:fubini-pruning}
 2|\cK_n|L_n\tau_n
 \le 2^{1+Cn-.01N}
 \le 2^{-.005N}
\end{equation}
for the chosen sufficiently large $M_0$ and all sufficiently
large $n$.

Every retained cell is its original unary cell and has mass
at least $\tau_n$.  We have not intersected the retained
sets over different opposite parameter values; such an
intersection could reduce a cell below its threshold.
Instead all subsequent estimates are made for this fixed
parameter pair and its unaltered cells, and then averaged.
This preserves both their mass bounds and the product of
the two conditional orientation laws.

\paragraph{Image entropy inside a retained cell.}
By \Cref{lem:small-tables}, the $k$ nominal key directions
are jointly independent modulo the pins, with rank computed
separately by component and sign and jointly across the two
endpoints.  Suppose a further tuple has rank $t\ge1$ modulo
pins and these key directions.  Apply
\eqref{eq:final-leaf-minentropy} to the combined tuple of
rank $k+t$, prescribe the known key images and any images
of the further tuple, and divide by the cell mass at least
$\tau_n$.  Its largest conditional point mass is at most
\begin{align}\label{eq:cell-pmax}
 2^{-(1-2\zeta)(k+t)N+(k+.01)N}
 &=2^{[-(1-2\zeta)t+.01+2\zeta k]N}
 \le2^{-.95tN}.
\end{align}
Here $k\le k_{\max}$ and
$\zeta=1/(1000k_{\max})$, so $.01+2\zeta k\le .012$.
The inequality holds for every $t\ge1$, including tuples
whose rank grows with $n$.  Further restrictions defining
the cell cause no additional loss: the numerator was bounded
by an event containing the entire cell and the prescribed
images, and only the original cell mass has been divided out.

\paragraph{Testing the remaining scalar Gram entries.}
For a retained cell pair the abstract target is fixed.  Let
$s_n$ be the number of its tested primal-channel entries in
the two cross Gram blocks.  By
\Cref{thm:gradient-realization} and the final choice of $M_0$,
\[
 s_n\le16|\cE|h\dim\cB<.01N.
\]
Expand the indicator that all these entries agree with the
target as a sum of $2^{s_n}$ binary characters, with the
factor $2^{-s_n}$.

A character has a coefficient tensor in the direct sum of
the two cross nominal tensor products, over all components.
Take its image after quotienting every nominal factor by
pins and keys.  If this image is zero, its value depends
only on frozen cross entries.  More explicitly the kernel
of a tensor-product quotient is the sum of tensors with a
frozen factor on at least one side; all their pairings have
been recorded.  The actual and target Grams agree there,
so the disagreement character is identically one.

Otherwise let $t\ge1$ be the sum of the ranks of the quotient
coefficient matrices, over components and the two cross
orientations.  A rank factorization expresses its phase,
up to a constant and two separate phases involving frozen
images, as
\[
 \sum_{a=1}^t X_a\cdot Y_a.
\]
On each unit the corresponding coefficient directions have
joint rank $t$ modulo its frozen space.  The two different
cross orientations use opposite nominal signs, so their
ranks add.  By \eqref{eq:cell-pmax}, each tuple of actual
images has largest point mass at most $2^{-.95tN}$ under
its conditional cell law.  The two laws are independent.
Absorb the separated phases into bounded weights.
\Cref{lem:walsh}, in dimension $tN$, bounds the absolute
character mean by
\begin{equation}\label{eq:gram-character-bound}
 2^{tN/2}\bigl(2^{-.95tN}2^{-.95tN}\bigr)^{1/2}
 =2^{-.45tN}\le2^{-.45N}.
\end{equation}

There is at least the trivial character with value one, and
all other zero-quotient characters also have value one.
Consequently the conditional agreement probability on every
retained cell pair is at least
\[
 2^{-s_n}\bigl(1-2^{s_n}2^{-.45N}\bigr)
 \ge2^{-s_n-1}\ge2^{-.02N}
\]
for all sufficiently large $n$.  No independence of the
individual tested Gram bits is asserted or required.

\paragraph{Averaging and the exponent.}
The assumed overlap is subexponentially large.  The error
in \eqref{eq:fubini-pruning} is exponentially small, so the
retained collision probability, after averaging leaves and
independent parameters, is at least
$2^{-o(N)}/|\cK_n|$.  Integrating the preceding scalar
agreement bound gives four holes with probability at least
\[
 2^{-(k+.02)N-o(N)}
\]
under $\widetilde\sigma^2$.  This implication is valid even
if many successful queries certify the same orientation
pair: the fresh query experiment is a probability space,
and its successful event is contained in the event that
the underlying orientations have all four holes.
The original law dominates its restricted normalized law
by the restriction mass; for two independent draws this
costs only another $2^{-o(N)}$.

If there are boundedly many options, sum the corresponding
overlap quantities and choose one whose expectation is at
least the sum divided by their number, or partition into
these fixed cases.  Shapes of $\cK_n$ and varying finite
numerical formats can likewise be separated into boundedly
many cases.  These constant factors are absorbed by the
slack from $.02$ to $.03$ in
\eqref{eq:collision-conclusion}.  Finally
\[
 k_{\max}+.03=56(g-1)+.03<100g,
\]
so \eqref{eq:collision-conclusion} contradicts a sequence
whose four-hole probabilities are less than $2^{-100gN}$.
\end{proof}
